% Options for packages loaded elsewhere
\PassOptionsToPackage{unicode}{hyperref}
\PassOptionsToPackage{hyphens}{url}
\PassOptionsToPackage{dvipsnames,svgnames,x11names}{xcolor}
%
\documentclass[
  11pt,
  letterpaper,
]{article}
\usepackage{amsmath,amssymb}
\usepackage{iftex}
\ifPDFTeX
  \usepackage[T1]{fontenc}
  \usepackage[utf8]{inputenc}
  \usepackage{textcomp} % provide euro and other symbols
\else % if luatex or xetex
  \usepackage{unicode-math} % this also loads fontspec
  \defaultfontfeatures{Scale=MatchLowercase}
  \defaultfontfeatures[\rmfamily]{Ligatures=TeX,Scale=1}
\fi
\usepackage{lmodern}
\ifPDFTeX\else
  % xetex/luatex font selection
\fi
% Use upquote if available, for straight quotes in verbatim environments
\IfFileExists{upquote.sty}{\usepackage{upquote}}{}
\IfFileExists{microtype.sty}{% use microtype if available
  \usepackage[]{microtype}
  \UseMicrotypeSet[protrusion]{basicmath} % disable protrusion for tt fonts
}{}
\makeatletter
\@ifundefined{KOMAClassName}{% if non-KOMA class
  \IfFileExists{parskip.sty}{%
    \usepackage{parskip}
  }{% else
    \setlength{\parindent}{0pt}
    \setlength{\parskip}{6pt plus 2pt minus 1pt}}
}{% if KOMA class
  \KOMAoptions{parskip=half}}
\makeatother
\usepackage{xcolor}
\usepackage{longtable,booktabs,array}
\usepackage{calc} % for calculating minipage widths
% Correct order of tables after \paragraph or \subparagraph
\usepackage{etoolbox}
\makeatletter
\patchcmd\longtable{\par}{\if@noskipsec\mbox{}\fi\par}{}{}
\makeatother
% Allow footnotes in longtable head/foot
\IfFileExists{footnotehyper.sty}{\usepackage{footnotehyper}}{\usepackage{footnote}}
\makesavenoteenv{longtable}
\usepackage{graphicx}
\makeatletter
\def\maxwidth{\ifdim\Gin@nat@width>\linewidth\linewidth\else\Gin@nat@width\fi}
\def\maxheight{\ifdim\Gin@nat@height>\textheight\textheight\else\Gin@nat@height\fi}
\makeatother
% Scale images if necessary, so that they will not overflow the page
% margins by default, and it is still possible to overwrite the defaults
% using explicit options in \includegraphics[width, height, ...]{}
\setkeys{Gin}{width=\maxwidth,height=\maxheight,keepaspectratio}
% Set default figure placement to htbp
\makeatletter
\def\fps@figure{htbp}
\makeatother
\setlength{\emergencystretch}{3em} % prevent overfull lines
\providecommand{\tightlist}{%
  \setlength{\itemsep}{0pt}\setlength{\parskip}{0pt}}
\setcounter{secnumdepth}{-\maxdimen} % remove section numbering
\usepackage[utf8]{inputenc}
\usepackage[T1]{fontenc}
\usepackage{lmodern}
\usepackage{textcomp}
\usepackage{pifont}
\usepackage{microtype}
\usepackage[margin=1in]{geometry}
\usepackage{longtable}
\usepackage{booktabs}
\usepackage{array}
\usepackage{xurl}
\setlength{\emergencystretch}{3em}
\setlength{\tabcolsep}{3pt}
% Allow line breaks after underscores (e.g. "sensor_log") so long
% identifiers wrap instead of overflowing narrow table columns.
\makeatletter
\let\bpd@origunderscore\_
\renewcommand{\_}{\bpd@origunderscore\allowbreak}
\makeatother
\ifLuaTeX
  \usepackage{selnolig}  % disable illegal ligatures
\fi
\IfFileExists{bookmark.sty}{\usepackage{bookmark}}{\usepackage{hyperref}}
\IfFileExists{xurl.sty}{\usepackage{xurl}}{} % add URL line breaks if available
\urlstyle{same}
\hypersetup{
  pdftitle={Entropy-Minimized Byte-Plane Decomposition: Auto-Detection and the Universality Gap in Lossless Compression},
  pdfauthor={Steffen Goehring},
  colorlinks=true,
  linkcolor={NavyBlue},
  filecolor={Maroon},
  citecolor={NavyBlue},
  urlcolor={NavyBlue},
  pdfcreator={LaTeX via pandoc}}

\title{Entropy-Minimized Byte-Plane Decomposition: Auto-Detection and
the Universality Gap in Lossless Compression}
\author{Steffen Goehring}
\date{August 2026}

\begin{document}
\maketitle

\hypertarget{abstract}{%
\subsection{Abstract}\label{abstract}}

Byte-plane decomposition --- reordering serialized multi-byte data by
byte position before compression --- is widely deployed in scientific
data pipelines (HDF5, Blosc, Apache Parquet) but requires the user or
application schema to specify the field width. We present two
\textbf{schema-free algorithms for automatically detecting the optimal
decomposition width} by minimizing average per-plane Shannon entropy:
one over a fixed candidate set, and one using sampled byte-match
frequency to detect arbitrary widths without a predefined set. The
algorithms are validated on data with record widths from 2 to 36 bytes,
including the SAO Star Catalog where width=28 is correctly identified
while all power-of-two widths fail. We additionally provide the first
\textbf{systematic characterization of the universality gap} between
structure-aware and structure-blind lossless compression. Across five
datasets --- synthetic sensor telemetry, real-world IoT sensors, MRI
medical imaging, astronomical catalog data, and financial order book
data --- the gap ranges from 1.00× to 40.8× and correlates with the
MSB-to-LSB entropy gradient and absolute MSB entropy within multi-byte
fields. We characterize the three-class byte-plane partition (constant,
sparse, medium-entropy) that emerges after decomposition and delta
encoding, and show it follows from the autoregressive structure of
typical binary field sequences. The synthetic sensor telemetry result
(40.8× gap) represents an upper bound on structured data with
near-constant MSBs; the real-world results (1.00×--1.19×) show the range
for data with higher per-field entropy. On the Intel Berkeley Lab sensor
dataset (2.2 million real IoT readings), the algorithm auto-detects
width=24 and achieves a 19\% compression improvement over ZSTD. On
NAS\-DAQ ITCH financial data (3.7 million order records), Algorithm 2
correctly detects the 36-byte record width, but the per-block
compression comparison confirms that decomposition cannot help (gap =
1.00×) --- demonstrating that entropy reduction alone does not guarantee
compression improvement. All results are validated on the BPD
implementation, a block-independent ZSTD-backed compression platform
implementing schema-free width detection and per-block strategy
selection.

\begin{center}\rule{0.5\linewidth}{0.5pt}\end{center}

\hypertarget{introduction}{%
\subsection{1. Introduction}\label{introduction}}

Lossless compression of binary data is dominated by universal
compressors --- algorithms like ZSTD {[}1{]}, LZ4, and zlib that operate
on raw byte streams without knowledge of the data's structure. These
compressors are optimal in the sense of Lempel-Ziv theory: they
asymptotically achieve the entropy rate of any stationary ergodic source
{[}2{]}. However, real binary data is often \textbf{not} a stationary
byte stream. Serialized structs, sensor telemetry, database pages, and
scientific data formats produce byte sequences where multi-byte integers
are interleaved --- a pattern that bounded-memory LZ compressors cannot
efficiently model when the interleaving period exceeds their effective
context window.

Byte-plane decomposition addresses this by rearranging an n-byte block
with w-byte records into w separate planes, each containing all bytes at
the same position within the record. This technique appears in HDF5's
shuffle filter {[}3{]}, Blosc's bitshuffle/bytedelta {[}4, 5{]}, and
Apache Parquet's BYTE\_STREAM\_SPLIT encoding {[}6{]}. Despite
widespread deployment, \textbf{no published algorithm automatically
detects the optimal decomposition width}, and \textbf{no formal analysis
quantifies the compression improvement} as a function of data
characteristics.

We make two contributions:\footnote{Implementation and reproducibility
  scripts are available at
  \url{https://github.com/leeroy37/byte-plane-decomposition}.}

\begin{enumerate}
\def\labelenumi{\arabic{enumi}.}
\item
  \textbf{Two entropy-minimized width detection algorithms:} Algorithm 1
  searches a fixed candidate set; Algorithm 2 uses sampled byte-match
  frequency to detect periodic byte repetitions at arbitrary lags,
  eliminating the predefined candidate set. Both find the optimal
  decomposition width w* by minimizing average per-plane Shannon
  entropy. We formalize the algorithms, analyze correctness and failure
  modes, and validate on data with struct widths from 2 to 36 bytes.
\item
  \textbf{An empirical characterization of the universality gap} between
  structure-aware (decomposition + ZSTD) and structure-blind (plain
  ZSTD) compression across five datasets (one synthetic, four
  real-world), with a theoretical explanation based on the MSB-to-LSB
  entropy gradient within multi-byte fields.
\end{enumerate}

\begin{center}\rule{0.5\linewidth}{0.5pt}\end{center}

\hypertarget{background}{%
\subsection{2. Background}\label{background}}

\hypertarget{byte-plane-decomposition}{%
\subsubsection{2.1 Byte-Plane
Decomposition}\label{byte-plane-decomposition}}

Given an input block of n bytes interpreted as n/w records of width w,
byte-plane decomposition produces w planes of size n/w each, where plane
p contains byte position p from all records:

\[\text{plane}_p[i] = \text{input}[i \cdot w + p], \quad 0 \le p < w, \quad 0 \le i < n/w\]

This is a deterministic permutation --- total information is preserved.
The benefit comes from grouping bytes with similar statistical
properties, enabling more efficient entropy coding by the downstream
compressor.

\textbf{Delta encoding} is typically applied per-plane after
decomposition: \(\text{delta}[i] = \text{plane}[i] - \text{plane}[i-1]\)
(with wraparound arithmetic). For slowly-changing fields, delta encoding
reduces most values to near-zero, further concentrating the probability
mass.

\hypertarget{prior-art}{%
\subsubsection{2.2 Prior Art}\label{prior-art}}

\begin{longtable}[]{@{}
  >{\raggedright\arraybackslash}p{(\columnwidth - 4\tabcolsep) * \real{0.1860}}
  >{\raggedright\arraybackslash}p{(\columnwidth - 4\tabcolsep) * \real{0.4419}}
  >{\raggedright\arraybackslash}p{(\columnwidth - 4\tabcolsep) * \real{0.3721}}@{}}
\toprule\noalign{}
\begin{minipage}[b]{\linewidth}\raggedright
System
\end{minipage} & \begin{minipage}[b]{\linewidth}\raggedright
Width Specification
\end{minipage} & \begin{minipage}[b]{\linewidth}\raggedright
Auto-Detection
\end{minipage} \\
\midrule\noalign{}
\endhead
\bottomrule\noalign{}
\endlastfoot
HDF5 shuffle filter {[}3{]} & User sets datatype size & No \\
Blosc bitshuffle {[}4{]} & User sets element size & No \\
Blosc bytedelta {[}5{]} & User sets typesize & No \\
Apache Parquet BYTE\_STREAM\_SPLIT {[}6{]} & Schema-derived & No
(requires schema) \\
TDT {[}7{]} & Entropy-based clustering & Partial (clusters byte
positions, doesn't detect width) \\
ZipNN {[}8{]} & AI model format-specific & No (hardcoded for
float16/32) \\
\textbf{This work} & \textbf{Entropy minimization (fixed set or
byte-match frequency)} & \textbf{Yes} \\
\end{longtable}

No prior system automatically detects the decomposition width from the
data alone.

\hypertarget{lz-compressor-limitations}{%
\subsubsection{2.3 LZ Compressor
Limitations}\label{lz-compressor-limitations}}

ZSTD and other LZ-family compressors find repeated byte sequences within
a sliding window and encode them as back-references. For interleaved
struct data with period \(w\), LZ matches shorter than \(w\) bytes
cannot capture the repeating field structure --- the compressor sees
each field's bytes interleaved with bytes from other fields, destroying
the local correlation patterns that LZ matching exploits.

Ferragina, Nitto, and Venturini {[}9{]} formalized this limitation: LZ77
is 8-optimal with respect to the zeroth-order empirical entropy \(H_0\)
but cannot be \(\lambda\)-optimal with respect to \(H_k\) for
\(k \ge 1\). The practical consequence is that the effective coding rate
of an LZ compressor on interleaved data approaches the \(H_0\) of the
mixed byte stream, which is substantially higher than the sum of
per-plane \(H_0\) values achievable after decomposition.

\begin{center}\rule{0.5\linewidth}{0.5pt}\end{center}

\hypertarget{entropy-minimized-width-detection}{%
\subsection{3. Entropy-Minimized Width
Detection}\label{entropy-minimized-width-detection}}

\hypertarget{algorithm}{%
\subsubsection{3.1 Algorithm}\label{algorithm}}

Given a block of n bytes, we seek the decomposition width \(w^*\) that
minimizes the average per-plane Shannon entropy:

\[w^* = \underset{w \in W}{\arg\min} \; \frac{1}{w} \sum_{p=0}^{w-1} H(\text{plane}_p^{(w)})\]

where \(W\) is the set of candidate widths and \(H(\cdot)\) denotes the
Shannon entropy of the byte frequency distribution of a plane.

\textbf{Algorithm 1: Entropy-Minimized Width Detection}

\begin{verbatim}
Input: byte block B of length n, candidate set W, threshold tau = 0.80
Output: optimal width w* or 0 (no decomposition)

H_base ← Shannon_entropy(B)
if H_base < 1.0 then return 0     // already low entropy
w* ← 0
H_best ← H_base

for each w in W do
    if n < 4w then continue        // need >=4 elements
    planes ← byte_shuffle(B, w)
    H_avg ← (1/w) sum_{p=0}^{w-1} Shannon_entropy(planes[p])
    if H_avg < tau · H_best then
        H_best ← H_avg
        w* ← w
return w*
\end{verbatim}

The threshold \(\tau = 0.80\) requires at least 20\% entropy reduction
to accept a width, preventing false positives from noise.

\hypertarget{candidate-width-selection-algorithm-1}{%
\subsubsection{3.2 Candidate Width Selection (Algorithm
1)}\label{candidate-width-selection-algorithm-1}}

Algorithm 1 uses a fixed candidate set
\(W = \{2, 3, 4, 5, 6, 7, 8, 10, 12, 14, 16, 20, 24, 28, 32\}\),
covering common struct sizes in systems programming, scientific data,
and network protocols. The computational cost is \(O(|W| \cdot n)\) ---
one shuffle and entropy computation per candidate width per block. This
set does not cover all possible widths; records wider than 32 bytes or
with widths not in \(W\) will not be detected.

\hypertarget{byte-match-based-candidate-detection-algorithm-2}{%
\subsubsection{3.2.1 Byte-Match-Based Candidate Detection (Algorithm
2)}\label{byte-match-based-candidate-detection-algorithm-2}}

To remove the fixed candidate set limitation, we introduce a second
algorithm that uses sampled byte-match frequency to detect periodic byte
repetitions at arbitrary lags, then validates only the top candidates
via entropy minimization.

\textbf{Algorithm 2: Byte-Match-Based Width Detection}

\begin{verbatim}
Input: byte block B of length n, threshold tau = 0.80, max lag L = 64
Output: optimal width w* or 0 (no decomposition)

H_base ← Shannon_entropy(B)
if H_base < 1.0 then return 0

// Phase 1: Sampled byte-match frequency
step ← max(1, n / 4096)              // sample ~4096 positions
for i = 0 to n-L by step do
    for lag = 2 to L do
        if B[i] = B[i + lag] then
            bmf[lag] ← bmf[lag] + 1
bmf[lag] ← bmf[lag] / sample_count   // normalize to [0, 1]

// Phase 2: Peak detection
threshold ← max(0.05, 0.5 · max(bmf))
C ← top 5 local maxima of bmf above threshold, sorted by bmf descending

// Phase 3: Add fallback widths
C ← C U {2, 4, 8}

// Phase 4: Entropy validation
for each w in C do
    if n < 4w then continue
    planes ← byte_shuffle(B, w)
    H_avg ← (1/w) sum_{p=0}^{w-1} Shannon_entropy(planes[p])
    if H_avg < tau · H_base then
        record (w, H_avg) as passing

// Phase 5: Harmonic suppression
if multiple widths pass then
    select w* with lowest H_avg
    for each smaller passing width w' that divides w*:
        if H_avg(w') < 1.10 · H_avg(w*) then w* ← w'
return w*
\end{verbatim}

\textbf{Byte-match frequency}
\(\text{bmf}(\ell) = P(B[i] = B[i + \ell])\) measures the fraction of
sampled positions where the byte value at position \(i\) equals the byte
value at position \(i + \ell\). This is not standard autocorrelation
\(R(\ell) = \sum (x_i - \bar{x})(x_{i+\ell} - \bar{x})\), but a simpler
indicator-function metric that is sufficient for period detection in
byte-valued data: structured data with record width \(w\) produces
byte-match peaks at lag \(w\) and its multiples \(2w, 3w, \ldots\)
because byte position \(p\) repeats every \(w\) bytes. By sampling
\textasciitilde4096 positions (rather than all \(n\)), the byte-match
phase costs \(O(L \cdot 4096) \approx 260K\) comparisons, negligible
relative to the \(O(n)\) shuffle cost per candidate (Section 7.5).

\textbf{Note on threshold behavior.} Algorithm 1 uses a cascading
threshold: \(H_{\text{best}}\) is updated as better widths are found, so
later candidates must beat \(\tau \times H_{\text{best\_so\_far}}\),
making the result order-dependent within the candidate set. Algorithm 2
compares all candidates against the fixed baseline
\(\tau \times H_{\text{base}}\), making it order-independent. On four of
five datasets the results are identical; NAS\-DAQ ITCH differs because
Algorithm 2 discovers width=36 (outside Algorithm 1's candidate set). In
general the cascading threshold is stricter and could produce different
results for marginal cases.

\textbf{Harmonic suppression} (Phase 5) addresses an artifact where
harmonics of the true period (e.g., \(2w\), \(3w\)) may also pass
entropy validation, because decomposition at width \(kw\) produces
\(kw\) planes that are still partially aligned with the underlying
\(w\)-byte structure. The byte-match frequency at the fundamental \(w\)
and its harmonics is theoretically equal for perfectly periodic data;
empirical differences arise from sampling noise or the peak detection
threshold. Regardless, preferring the smallest passing divisor selects
the fundamental period. The 10\% tolerance threshold is empirically
chosen to allow harmonics with marginally lower entropy while still
preferring the fundamental.

\textbf{Advantages over Algorithm 1:} (1) Supports arbitrary widths up
to \(L\) without a predefined candidate set --- width 36 (NAS\-DAQ ITCH),
width 28 (SAO catalog), width 24 (Intel Lab), and any non-standard width
are discovered automatically. For widths present in the fallback set (2,
4, 8), the byte-match phase may not produce a peak at the correct lag;
the fallback ensures these common widths are always tested. The
byte-match mechanism's primary contribution is discovering non-standard
widths. (2) Evaluates only 5--8 candidates per block instead of 15,
reducing the number of shuffle+entropy operations. (3) The fallback set
\(\{2, 4, 8\}\) ensures common power-of-two widths are always tested
even when byte-match sampling is insufficient (e.g., on short blocks).

\hypertarget{correctness-analysis}{%
\subsubsection{3.3 Correctness Analysis}\label{correctness-analysis}}

\textbf{When the algorithm succeeds:} If the data consists of repeated
w-byte records where different byte positions have different entropy
characteristics (e.g., MSBs are constant, LSBs vary), the correct width
w produces planes with the lowest average entropy. Any incorrect width
\(w' \ne w\) (and \(w' \nmid w\)) mixes bytes from different positions,
averaging their entropy upward.

\textbf{When the algorithm fails:} (1) If all byte positions within the
record have similar entropy (e.g., all bytes of a uniformly distributed
32-bit float), no width reduces average entropy. (2) If the record width
is not in the candidate set \(W\), it will be missed (Algorithm 1 only;
Algorithm 2 detects arbitrary widths up to \(L\)). (3) If the block
contains mixed record types, the dominant type determines the detected
width. (4) Variable-length serialization formats (protobuf, MessagePack,
CBOR) are correctly rejected because no fixed period produces consistent
entropy reduction.

\hypertarget{validation}{%
\subsubsection{3.4 Validation}\label{validation}}

\textbf{SAO Star Catalog (28-byte records):} The algorithm correctly
identifies w* = 28 from the candidate set. All smaller widths (2, 4, 8,
14, 16) fail the 20\% threshold --- they produce average plane entropies
of 6.19--7.17 bits, reductions of only 3--16\% versus the baseline of
7.40 bits. Width 28 produces average plane entropy of 5.39 bits, a 27\%
reduction. Width 56 (2 × 28) also passes, confirming the periodic
structure.

\begin{longtable}[]{@{}llll@{}}
\toprule\noalign{}
Width & Avg Plane Entropy & Entropy Reduction & Accepted \\
\midrule\noalign{}
\endhead
\bottomrule\noalign{}
\endlastfoot
2 & 7.17 & 3\% & No (\textless{} 20\%) \\
4 & 6.82 & 8\% & No (\textless{} 20\%) \\
8 & 6.80 & 8\% & No (\textless{} 20\%) \\
14 & 6.19 & 16\% & No (\textless{} 20\%) \\
\textbf{28} & \textbf{5.39} & \textbf{27\%} & \textbf{Yes} \\
32 & 6.73 & 9\% & No (\textless{} 20\%) \\
\end{longtable}

\textbf{Sensor telemetry (8-byte records):} The algorithm identifies w*
= 8. Average plane entropy drops from 5.4 bits to 2.8 bits (49\%
reduction), well above the 20\% threshold.

\hypertarget{threshold-sensitivity}{%
\subsubsection{3.5 Threshold Sensitivity}\label{threshold-sensitivity}}

The threshold \(\tau = 0.80\) was chosen to balance detection
sensitivity against false positives. We sweep \(\tau\) from 0.65 to 0.95
across all datasets:

\begin{longtable}[]{@{}
  >{\raggedright\arraybackslash}p{(\columnwidth - 14\tabcolsep) * \real{0.0889}}
  >{\raggedright\arraybackslash}p{(\columnwidth - 14\tabcolsep) * \real{0.1444}}
  >{\raggedright\arraybackslash}p{(\columnwidth - 14\tabcolsep) * \real{0.1444}}
  >{\raggedright\arraybackslash}p{(\columnwidth - 14\tabcolsep) * \real{0.1000}}
  >{\raggedright\arraybackslash}p{(\columnwidth - 14\tabcolsep) * \real{0.1222}}
  >{\raggedright\arraybackslash}p{(\columnwidth - 14\tabcolsep) * \real{0.1444}}
  >{\raggedright\arraybackslash}p{(\columnwidth - 14\tabcolsep) * \real{0.0889}}
  >{\raggedright\arraybackslash}p{(\columnwidth - 14\tabcolsep) * \real{0.1667}}@{}}
\toprule\noalign{}
\begin{minipage}[b]{\linewidth}\raggedright
\(\tau\)
\end{minipage} & \begin{minipage}[b]{\linewidth}\raggedright
sensor (w=8)
\end{minipage} & \begin{minipage}[b]{\linewidth}\raggedright
intel (w=24)
\end{minipage} & \begin{minipage}[b]{\linewidth}\raggedright
mr (w=2)
\end{minipage} & \begin{minipage}[b]{\linewidth}\raggedright
sao (w=28)
\end{minipage} & \begin{minipage}[b]{\linewidth}\raggedright
dickens (w=0)
\end{minipage} & \begin{minipage}[b]{\linewidth}\raggedright
NAS\-DAQ
\end{minipage} & \begin{minipage}[b]{\linewidth}\raggedright
Correct (det.)
\end{minipage} \\
\midrule\noalign{}
\endhead
\bottomrule\noalign{}
\endlastfoot
0.65 & 8 \ding{51} & 24 \ding{51} & --- \ding{51} & --- \ding{55} & ---
\ding{51} & --- & 4/5 \\
0.70 & 8 \ding{51} & 12 \ding{55} & --- \ding{51} & --- \ding{55} & ---
\ding{51} & --- & 3/5 \\
0.75 & 16 \ding{55} & 24 \ding{51} & --- \ding{51} & 28 \ding{51} & ---
\ding{51} & 12 & 4/5 \\
\textbf{0.80} & \textbf{8 \ding{51}} & \textbf{24 \ding{51}} & \textbf{2
\ding{51}} & \textbf{28 \ding{51}} & \textbf{--- \ding{51}} &
\textbf{12} & \textbf{5/5} \\
0.85 & 8 \ding{51} & 24 \ding{51} & 20 \ding{55} & 14 \ding{55} & ---
\ding{51} & 12 & 3/5 \\
0.90 & 8 \ding{51} & 24 \ding{51} & 2 \ding{51} & 28 \ding{51} & ---
\ding{51} & 12 & 5/5 \\
0.95 & 32 \ding{55} & 24 \ding{51} & 2 \ding{51} & 28 \ding{51} & ---
\ding{51} & 12 & 4/5 \\
\end{longtable}

NAS\-DAQ is excluded from the ``Correct'' count because its
detection-stage behavior requires special interpretation: Algorithm 1
detects width=12 (a divisor of the true width 36) at all thresholds
\(\tau \ge 0.75\), with 22--24\% entropy reduction. This is not a false
positive --- the 36-byte records genuinely contain 12-byte periodic
sub-structure --- but it does not produce a compression improvement
(\(G = 1.00\times\)). Algorithm 2 detects the true width=36 with 48\%
entropy reduction, also without compression improvement.

\(\tau = 0.80\) is the most conservative threshold achieving 5/5
correctness on the non-NAS\-DAQ datasets. \(\tau = 0.90\) also achieves
5/5, but sits between two failing thresholds (0.85 at 3/5, 0.95 at 4/5),
suggesting it may be a fragile coincidence on this dataset collection.
\(\tau = 0.80\) provides more margin below the SAO detection point (27\%
reduction) and is therefore preferred. Below 0.80, SAO is missed; at
0.85, both mr and sao produce incorrect widths.

The two-stage design (detection + per-block compression comparison)
provides robustness beyond the threshold alone: NAS\-DAQ's detected widths
are always rejected at the compression comparison stage because
shuffle+delta+ZSTD does not produce shorter output than plain ZSTD,
resulting in no regression at any threshold. This defense-in-depth means
the system never regresses even when the detection stage identifies
structure that does not benefit compression.

\begin{center}\rule{0.5\linewidth}{0.5pt}\end{center}

\hypertarget{the-universality-gap}{%
\subsection{4. The Universality Gap}\label{the-universality-gap}}

\hypertarget{definition}{%
\subsubsection{4.1 Definition}\label{definition}}

We define the \textbf{universality gap} \(G\) as the ratio of
compression achieved by structure-aware compression (decomposition +
delta + ZSTD) to structure-blind compression (plain ZSTD) on the same
data:

\[G = \frac{R_{\text{shuffle+delta+ZSTD}}}{R_{\text{ZSTD}}}\]

where \(R\) denotes the compression ratio (original size / compressed
size). \(G > 1\) indicates that byte-plane decomposition improves
compression; \(G \gg 1\) indicates a large gap that structure-blind
compressors cannot close.

\hypertarget{experimental-measurements}{%
\subsubsection{4.2 Experimental
Measurements}\label{experimental-measurements}}

We measured \(G\) on five datasets with distinct structural properties,
using 64 KB independent blocks and ZSTD level 3 as the baseline:

\begin{longtable}[]{@{}
  >{\raggedright\arraybackslash}p{(\columnwidth - 10\tabcolsep) * \real{0.1304}}
  >{\raggedright\arraybackslash}p{(\columnwidth - 10\tabcolsep) * \real{0.1159}}
  >{\raggedright\arraybackslash}p{(\columnwidth - 10\tabcolsep) * \real{0.1014}}
  >{\raggedright\arraybackslash}p{(\columnwidth - 10\tabcolsep) * \real{0.2754}}
  >{\raggedright\arraybackslash}p{(\columnwidth - 10\tabcolsep) * \real{0.3043}}
  >{\raggedright\arraybackslash}p{(\columnwidth - 10\tabcolsep) * \real{0.0725}}@{}}
\toprule\noalign{}
\begin{minipage}[b]{\linewidth}\raggedright
Dataset
\end{minipage} & \begin{minipage}[b]{\linewidth}\raggedright
Format
\end{minipage} & \begin{minipage}[b]{\linewidth}\raggedright
Width
\end{minipage} & \begin{minipage}[b]{\linewidth}\raggedright
\(R_{\text{ZSTD}}\)
\end{minipage} & \begin{minipage}[b]{\linewidth}\raggedright
\(R_{\text{shuffle}}\)
\end{minipage} & \begin{minipage}[b]{\linewidth}\raggedright
\(G\)
\end{minipage} \\
\midrule\noalign{}
\endhead
\bottomrule\noalign{}
\endlastfoot
Sensor telemetry (synth.) & 8B: uint32 ts + int16 temp + int16 hum & 8 &
4.25:1 & 173.5:1 & \textbf{40.8×} \\
Intel Lab sensors (real) & 24B: uint32 epoch + uint32 moteid + 4×float32
& 24 & 3.58:1 & 4.26:1 & \textbf{1.19×} \\
MRI brain scan (Silesia \texttt{mr}) & 16-bit DICOM pixels & 2 & 2.72:1
& 2.76:1 & \textbf{1.02×} \\
SAO Star Catalog (Silesia \texttt{sao}) & 28B astronomical records & 28
& 1.25:1 & 1.26:1 & \textbf{1.01×} \\
NAS\-DAQ ITCH Add Orders (real) & 36B: 6B nanos ts + 8B order ref + 8B
stock + \ldots{} & 36 & 2.80:1 & 2.80:1 & \textbf{1.00×} \\
\end{longtable}

*Algorithm 1 detects width=12 (a divisor of 36 present in the candidate
set) with 23\% entropy reduction. Algorithm 2 detects the true width=36
with 48\% entropy reduction. In both cases, the per-block compression
comparison rejects shuffle because shuffle+delta+ZSTD does not produce
shorter output than plain ZSTD on any block. The \(R_{\text{shuffle}}\)
column shows the result at w=36 (Algorithm 2); the result at w=12
(Algorithm 1) is identical (2.80:1).

\hypertarget{the-entropy-gradient-explanation}{%
\subsubsection{4.3 The Entropy Gradient
Explanation}\label{the-entropy-gradient-explanation}}

The universality gap correlates with the \textbf{MSB-to-LSB entropy
gradient} --- the difference in Shannon entropy between the most
significant and least significant byte planes of each multi-byte field:

\begin{longtable}[]{@{}lllll@{}}
\toprule\noalign{}
Dataset & MSB Entropy & LSB Entropy & Gradient & Gap \(G\) \\
\midrule\noalign{}
\endhead
\bottomrule\noalign{}
\endlastfoot
Sensor telemetry (synth.) & 0.00 bits & 5.00 bits & 5.00 & 40.8× \\
Intel Lab sensors (real) & 0.17 bits & 6.07 bits & 5.90 & 1.19× \\
MRI (Silesia \texttt{mr}) & 0.47 bits & 5.05 bits & 4.58 & 1.02× \\
SAO (Silesia \texttt{sao}) & 3.19 bits & 6.08 bits & 2.89 & 1.01× \\
NAS\-DAQ ITCH (real) & \textasciitilde6.0 bits & \textasciitilde7.5 bits &
\textasciitilde1.5 & 1.00× \\
\end{longtable}

\textbf{Observation.} The gradient alone does not predict the gap:
synthetic sensor data and Intel Lab sensors have similar gradients (5.00
vs 5.90) but differ by 34× in gap (40.8× vs 1.19×). The critical
additional factor is the \textbf{absolute MSB entropy}. When MSBs are
exactly 0.00 bits (synthetic: constant-increment timestamps), the
separated MSB planes compress to nearly zero bytes. When MSBs are 0.17
bits (Intel Lab: variable epochs, 54 sensor IDs), the separated planes
still have non-trivial content. Both the gradient and the absolute MSB
entropy are necessary predictors of the gap.

\hypertarget{qualitative-analysis}{%
\subsubsection{4.4 Qualitative Analysis}\label{qualitative-analysis}}

Our five data points (Figure 1) suggest that the universality gap \(G\)
is governed by two factors:

\begin{enumerate}
\def\labelenumi{\arabic{enumi}.}
\item
  \textbf{The MSB-to-LSB entropy gradient}
  \(\Delta H = H(\text{LSB plane}) - H(\text{MSB plane})\): a large
  gradient means decomposition separates near-zero-entropy planes from
  high-entropy planes, giving ZSTD dramatically better input on the
  separated planes.
\item
  \textbf{The absolute MSB entropy}: even with a large gradient, if the
  MSB planes retain significant entropy (SAO: 3.19 bits), no plane
  becomes trivially compressible and the gap stays small.
\end{enumerate}

The mechanism is not simply that decomposition reduces total entropy ---
it cannot, since it is an information-preserving permutation. Rather,
decomposition converts the problem from coding a single high-entropy
interleaved stream into coding multiple streams of varying entropy.
LZ-family compressors exploit low-entropy streams far more efficiently
than they exploit low-entropy \emph{subsequences} embedded within a
high-entropy stream, because their match-finding operates on contiguous
windows.

\begin{figure}
\centering
\includegraphics[width=0.75\textwidth,height=\textheight]{universality_gap.pdf}
\caption{Universality gap \(G\) versus absolute MSB entropy for all five
datasets: sensor\_log (0.00, 40.8), Intel Lab (0.17, 1.19), MRI (0.47,
1.02), SAO (3.19, 1.01), NAS\-DAQ (6.0, 1.00). The gap decreases toward
1.0 as MSB entropy increases, confirming that near-zero MSB entropy is
the primary driver of large gaps.}
\end{figure}

Tightening this qualitative relationship into a formal bound remains an
open problem. It requires bounding the gap between ZSTD's effective
per-byte coding rate on periodically interleaved data (which approaches
\(H_0\) of the mixed stream) and the sum of per-plane coding rates after
decomposition. The framework of Ferragina et al.~{[}9{]} provides the
necessary LZ optimality bounds, but connecting these to the specific
structure of byte-interleaved data requires further analysis.

\begin{center}\rule{0.5\linewidth}{0.5pt}\end{center}

\hypertarget{the-three-class-plane-partition}{%
\subsection{5. The Three-Class Plane
Partition}\label{the-three-class-plane-partition}}

\hypertarget{empirical-observation}{%
\subsubsection{5.1 Empirical Observation}\label{empirical-observation}}

After byte-plane decomposition at the correct width followed by delta
encoding, each byte plane consistently falls into one of three entropy
classes:

\begin{longtable}[]{@{}
  >{\raggedright\arraybackslash}p{(\columnwidth - 6\tabcolsep) * \real{0.1250}}
  >{\raggedright\arraybackslash}p{(\columnwidth - 6\tabcolsep) * \real{0.3393}}
  >{\raggedright\arraybackslash}p{(\columnwidth - 6\tabcolsep) * \real{0.1607}}
  >{\raggedright\arraybackslash}p{(\columnwidth - 6\tabcolsep) * \real{0.3750}}@{}}
\toprule\noalign{}
\begin{minipage}[b]{\linewidth}\raggedright
Class
\end{minipage} & \begin{minipage}[b]{\linewidth}\raggedright
Entropy after delta
\end{minipage} & \begin{minipage}[b]{\linewidth}\raggedright
Example
\end{minipage} & \begin{minipage}[b]{\linewidth}\raggedright
Compression strategy
\end{minipage} \\
\midrule\noalign{}
\endhead
\bottomrule\noalign{}
\endlastfoot
\textbf{Constant} & \textless{} 0.01 bits/byte & Timestamp MSBs, field
sign bytes & Generative: store pattern seed only \\
\textbf{Sparse} & 0.01--0.50 bits/byte & Slowly-varying LSBs, carry
bytes & ZSTD: near-zero compressed size \\
\textbf{Medium} & 0.50--8.00 bits/byte & High-entropy coordinate data,
pixel values & ZSTD: standard compression \\
\end{longtable}

\hypertarget{formal-statement}{%
\subsubsection{5.2 Formal Statement}\label{formal-statement}}

We state two propositions: the first for the static (single-value) case,
the second for the sequential (autoregressive) case that models typical
binary data.

\textbf{Proposition 1 (Static case).} Let \(X\) be a w-byte unsigned
integer with values in \([0, V]\) where \(V < 256^w\). After byte-plane
decomposition, at least \(w - \lceil \log_{256}(V+1) \rceil\) MSB
byte-planes have entropy 0, since byte position \(p\) is identically
zero when \(V < 256^p\).

\emph{This is elementary but establishes the baseline: bounded integers
have zero-entropy MSB planes proportional to the unused dynamic range.}

\textbf{Proposition 2 (Autoregressive case; stated without formal
proof).} Let \(\{X_i\}\) be a sequence of w-byte unsigned integers with
\(X_i = X_{i-1} + \delta + \epsilon_i\), where \(\delta\) is a constant
increment. In the deterministic case (\(\epsilon_i = 0\)), after
byte-plane decomposition followed by delta encoding, the planes
partition into three classes. With noise (\(|\epsilon_i| \ll 256^p\) for
some threshold position \(p\)), constant-class planes become
near-constant with deviation probability bounded by
\(P(|\epsilon_i| \ge 256^q)\) for position \(q\), which is negligible
under the stated assumption. The three classes are:

\begin{enumerate}
\def\labelenumi{\arabic{enumi}.}
\item
  \textbf{Constant class} (positions \(q > p\) where
  \(\delta < 256^{q}\)): After delta encoding, plane \(q\) consists of a
  single anchor byte followed by a constant value (typically 0), since
  consecutive values of \(\lfloor X_i / 256^q \rfloor\) are identical
  whenever the increment \(\delta\) does not span byte boundary
  \(256^q\). Entropy \(\approx 0\).
\item
  \textbf{Sparse class} (position \(p\) and nearby): Plane \(p\) after
  delta encoding contains the carry bits from lower-position overflows.
  The carry frequency at position \(p\) is
  \((\delta \bmod 256^{p+1}) / 256^{p+1}\), producing a sparse
  distribution concentrated on 1--3 distinct values. Entropy in the
  range {[}0.01, 0.50{]} bits/byte.
\item
  \textbf{Medium class} (positions \(q < p\)): Lower byte planes contain
  the residuals \(\epsilon_i \bmod 256^{q+1}\) plus the deterministic
  pattern \(\delta \bmod 256^{q+1}\). After delta encoding, these reduce
  to the noise terms \(\epsilon_i\), which retain entropy proportional
  to \(H(\epsilon_i)\).
\end{enumerate}

\emph{The threshold position \(p\) is determined by the magnitude of
\(\delta\): \(p = \lfloor \log_{256}(\delta) \rfloor\). For sensor
telemetry with \(\delta = 1000\),
\(p = \lfloor \log_{256}(1000) \rfloor = 1\), predicting constant planes
at positions 2+ and a sparse carry plane at position 1. This matches our
experimental observations (Section 5.3).}

\textbf{Carry frequency prediction.} At the boundary position, the
carry-propagation frequency from plane \(p-1\) to plane \(p\) is:

\[f_{\text{carry}} = \frac{\delta \bmod 256}{256}\]

For \(\delta = 1000\): \(f_{\text{carry}} = 232/256 \approx 0.906\),
predicting a bimodal delta distribution with 90.6\% of one value and
9.4\% of another --- exactly matching the observed distribution on plane
1 of the sensor telemetry data.

\textbf{Note on endianness.} The algorithm is endianness-agnostic: it
identifies low-entropy planes regardless of byte order. The labels
``MSB'' and ``LSB'' in the three-class partition refer to the byte
position's entropy role, not the endianness convention. Little-endian
data (sensor\_log: LSB at position 0) and big-endian data (SAO star
catalog) both exhibit the partition, with the low-entropy planes
appearing at different position indices.

\hypertarget{validation-on-sensor-telemetry}{%
\subsubsection{5.3 Validation on Sensor
Telemetry}\label{validation-on-sensor-telemetry}}

Per-plane entropy profile (8-byte struct, width=8, averaged over 16
blocks of 64 KB):

\begin{longtable}[]{@{}
  >{\raggedright\arraybackslash}p{(\columnwidth - 8\tabcolsep) * \real{0.1346}}
  >{\raggedright\arraybackslash}p{(\columnwidth - 8\tabcolsep) * \real{0.1346}}
  >{\raggedright\arraybackslash}p{(\columnwidth - 8\tabcolsep) * \real{0.3462}}
  >{\raggedright\arraybackslash}p{(\columnwidth - 8\tabcolsep) * \real{0.2500}}
  >{\raggedright\arraybackslash}p{(\columnwidth - 8\tabcolsep) * \real{0.1346}}@{}}
\toprule\noalign{}
\begin{minipage}[b]{\linewidth}\raggedright
Plane
\end{minipage} & \begin{minipage}[b]{\linewidth}\raggedright
Field
\end{minipage} & \begin{minipage}[b]{\linewidth}\raggedright
Avg Entropy (raw)
\end{minipage} & \begin{minipage}[b]{\linewidth}\raggedright
After Delta
\end{minipage} & \begin{minipage}[b]{\linewidth}\raggedright
Class
\end{minipage} \\
\midrule\noalign{}
\endhead
\bottomrule\noalign{}
\endlastfoot
0 & timestamp LSB & 5.000 & 0.002 & Constant
(\(\delta \bmod 256 = 232\)) \\
1 & timestamp byte 1 & 8.000 & 0.451 & Sparse (carry propagation) \\
2 & timestamp byte 2 & 6.972 & 0.116 & Sparse \\
3 & timestamp MSB & 0.187 & 0.002 & Constant \\
4 & temp LSB & 1.681 & 0.013 & Sparse \\
5 & temp MSB & 0.000 & 0.002 & Constant \\
6 & humidity LSB & 1.717 & 0.013 & Sparse \\
7 & humidity MSB & 0.000 & 0.002 & Constant \\
\end{longtable}

Four planes (0, 3, 5, 7) are constant-class after delta encoding. Four
planes (1, 2, 4, 6) are sparse, including the carry-propagation plane at
position 1 (entropy 0.451, near the class boundary). The carry plane's
bimodal distribution (90.6\% \texttt{0x04}, 9.4\% \texttt{0x03}) is
determined by the LSB overflow frequency:
\((1000 \bmod 256) / 256 = 232/256 \approx 0.906\).

\begin{center}\rule{0.5\linewidth}{0.5pt}\end{center}

\hypertarget{experimental-setup}{%
\subsection{6. Experimental Setup}\label{experimental-setup}}

\hypertarget{platform}{%
\subsubsection{6.1 Platform}\label{platform}}

All experiments use the experimental platform, a C++20 lossless
compression format with block-independent 64 KB blocks. Each block is
independently compressed using the best strategy: byte-plane
decomposition + delta + ZSTD, or plain ZSTD (fallback). The per-block
strategy selection follows the MDL principle: both candidates are
evaluated and the shortest output is kept.

ZSTD level 3 is used throughout. ZSTD context is reused across blocks to
amortize initialization cost.

\hypertarget{datasets}{%
\subsubsection{6.2 Datasets}\label{datasets}}

\begin{longtable}[]{@{}
  >{\raggedright\arraybackslash}p{(\columnwidth - 8\tabcolsep) * \real{0.1837}}
  >{\raggedright\arraybackslash}p{(\columnwidth - 8\tabcolsep) * \real{0.1224}}
  >{\raggedright\arraybackslash}p{(\columnwidth - 8\tabcolsep) * \real{0.1633}}
  >{\raggedright\arraybackslash}p{(\columnwidth - 8\tabcolsep) * \real{0.2653}}
  >{\raggedright\arraybackslash}p{(\columnwidth - 8\tabcolsep) * \real{0.2653}}@{}}
\toprule\noalign{}
\begin{minipage}[b]{\linewidth}\raggedright
Dataset
\end{minipage} & \begin{minipage}[b]{\linewidth}\raggedright
Size
\end{minipage} & \begin{minipage}[b]{\linewidth}\raggedright
Source
\end{minipage} & \begin{minipage}[b]{\linewidth}\raggedright
Record Width
\end{minipage} & \begin{minipage}[b]{\linewidth}\raggedright
Description
\end{minipage} \\
\midrule\noalign{}
\endhead
\bottomrule\noalign{}
\endlastfoot
sensor\_log & 16 MB & Synthetic & 8 bytes & IoT telemetry: uint32
timestamp (\(\Delta\)=1000) + int16 temp (\(\sigma\)=0.3) + int16
humidity (\(\sigma\)=0.3) \\
intel\_lab & 50.8 MB & Intel Berkeley Lab {[}16{]} & 24 bytes & Real
IoT: 2.2M readings from 54 sensors; uint32 epoch + uint32 moteid +
4×float32 \\
mr & 9.97 MB & Silesia Corpus {[}10{]} & 2 bytes & MRI brain scan,
16-bit grayscale DICOM pixel data \\
sao & 7.25 MB & Silesia Corpus {[}10{]} & 28 bytes & SAO Star Catalog,
fixed-width astronomical records \\
NAS\-DAQ ITCH & 131.8 MB & NAS\-DAQ {[}17{]} & 36 bytes & Add Order
messages: 6B nanosecond timestamp + 8B order ref + 8B stock symbol +
fields \\
\end{longtable}

\textbf{Note on synthetic data.} The sensor\_log dataset is synthetic,
designed to represent the class of IoT/SCADA/telemetry data with
constant-increment timestamps and slowly-drifting sensor values. The
40.8× universality gap represents an upper bound on highly regular
structured data. The Intel Berkeley Lab dataset {[}16{]} provides a
real-world IoT comparison (2.2 million readings from 54 sensors,
converted from CSV to 24-byte binary records), showing a gap of 1.19×
--- substantially less than the synthetic case due to irregular
timestamps, variable sensor IDs, and missing readings. The Silesia
results (1.01×--1.02×) represent real data with higher per-field
entropy. Additional real-world binary datasets with estimated gaps in
the 2×--10× range include USGS LiDAR LAS point clouds (30-byte records)
{[}18{]} and LBNL/ICSI enterprise packet captures (90-byte records)
{[}19{]}. NAS\-DAQ ITCH 5.0 {[}17{]} was tested and showed no benefit (gap
= 1.00×), validating the algorithm's correct rejection of high-entropy
structured data.

\hypertarget{methodology}{%
\subsubsection{6.3 Methodology}\label{methodology}}

For each dataset: (1) compute block-level Shannon entropy, (2) run width
detection algorithm with candidate set
\(W = \{2,3,4,5,6,7,8,10,12,14,16,20,24,28,32\}\), (3) at the detected
width, compute per-plane entropy before and after delta encoding, (4)
compress with shuffle+delta+ZSTD and compare against plain ZSTD on the
same block, (5) compute the universality gap \(G\).

All width detection results in Sections 7.1--7.4 use Algorithm 1.
Algorithm 2 produces identical width selections on four of five
datasets; on NAS\-DAQ ITCH, Algorithm 2 detects the true width=36 (outside
Algorithm 1's candidate set) while Algorithm 1 detects w=12. In both
cases the per-block compression comparison rejects the shuffle result.
Section 7.1 includes Algorithm 2's byte-match profiles for validation.

All entropy values are order-0 Shannon entropy of the byte frequency
distribution, measured in bits per byte (range {[}0, 8{]}). Width
detection (Algorithms 1 and 2) evaluates raw (pre-delta) per-plane
entropy. The three-class partition analysis (Section 5) uses post-delta
entropy.

\begin{center}\rule{0.5\linewidth}{0.5pt}\end{center}

\hypertarget{results}{%
\subsection{7. Results}\label{results}}

\hypertarget{width-detection-accuracy}{%
\subsubsection{7.1 Width Detection
Accuracy}\label{width-detection-accuracy}}

The algorithm correctly identifies the struct width for all four
structured datasets and correctly rejects decomposition on all
non-structured data:

\textbf{Successful detections:}

\begin{longtable}[]{@{}
  >{\raggedright\arraybackslash}p{(\columnwidth - 8\tabcolsep) * \real{0.1304}}
  >{\raggedright\arraybackslash}p{(\columnwidth - 8\tabcolsep) * \real{0.1594}}
  >{\raggedright\arraybackslash}p{(\columnwidth - 8\tabcolsep) * \real{0.2174}}
  >{\raggedright\arraybackslash}p{(\columnwidth - 8\tabcolsep) * \real{0.2754}}
  >{\raggedright\arraybackslash}p{(\columnwidth - 8\tabcolsep) * \real{0.2174}}@{}}
\toprule\noalign{}
\begin{minipage}[b]{\linewidth}\raggedright
Dataset
\end{minipage} & \begin{minipage}[b]{\linewidth}\raggedright
True Width
\end{minipage} & \begin{minipage}[b]{\linewidth}\raggedright
Detected Width
\end{minipage} & \begin{minipage}[b]{\linewidth}\raggedright
Entropy Reduction
\end{minipage} & \begin{minipage}[b]{\linewidth}\raggedright
Detection Time
\end{minipage} \\
\midrule\noalign{}
\endhead
\bottomrule\noalign{}
\endlastfoot
sensor\_log & 8 & 8 & 49\% (5.4 → 2.8 bits) & \textasciitilde100
\(\mu\)s/block \\
intel\_lab & 24 & 24 & 50\% (6.1 → 3.0 bits) & \textasciitilde150
\(\mu\)s/block \\
mr & 2 & 2 & 20--22\% on activating blocks & \textasciitilde100
\(\mu\)s/block \\
sao & 28 & 28 & 27\% (7.4 → 5.4 bits) & \textasciitilde200
\(\mu\)s/block \\
NAS\-DAQ ITCH & 36 & 12 (Alg 1) / 36 (Alg 2) & 23\% at w=12, 48\% at w=36
& \textasciitilde200 \(\mu\)s/block \\
\end{longtable}

Note on mr: Width=2 is detected on approximately 30\% of blocks (those
with low-entropy uniform-tissue regions, block entropy
\textasciitilde2.2--2.6). On these blocks, the per-plane entropy
reduction is 20--22\%, meeting the \(\tau\)=0.80 threshold. On the
remaining 70\% of blocks (higher-entropy regions, block entropy
\textgreater3.0), no width passes the threshold and the algorithm
correctly falls back to plain ZSTD. The aggregate entropy reduction
across all blocks is \textasciitilde14\%.

Note on NAS\-DAQ ITCH: Algorithm 1 detects width=12 (a divisor of 36,
present in the candidate set) with 23\% entropy reduction. Algorithm 2
detects the true record width of 36 with 48\% entropy reduction --- a
strong byte-match peak at lag=36 (frequency 0.84). However, in both
cases the per-block compression comparison rejects the shuffle result:
shuffle+delta+ZSTD does not compress smaller than plain ZSTD on any
block, yielding \(G = 1.00\times\). This demonstrates that entropy
reduction (even 48\%) does not guarantee compression improvement ---
ZSTD's LZ matching already captures the interleaved patterns
effectively. The two-stage design (detection + per-block comparison)
correctly prevents regression.

\textbf{Correct rejections (Silesia Corpus non-structured files):}

\begin{longtable}[]{@{}lllll@{}}
\toprule\noalign{}
File & Type & Block Entropy & Detected Width & Activated? \\
\midrule\noalign{}
\endhead
\bottomrule\noalign{}
\endlastfoot
dickens & English text & 4.61 & 0 & No \\
mozilla & x86 executable & 4.06 & 0 & No \\
nci & Chemical DB (text) & 2.42 & 0 & No \\
ooffice & Windows DLL & 6.29 & 0 & No \\
osdb & Synthetic DB & 6.59 & 0 & No \\
reymont & Polish text (PDF) & 4.84 & 0 & No \\
samba & Source code & 4.89 & 0 & No \\
webster & HTML dictionary & 5.04 & 0 & No \\
xml & XML text & 5.23 & 0 & No \\
\end{longtable}

The algorithm produces \textbf{zero false positives} on the entire
Silesia Corpus text and executable files. No candidate width achieves
the 20\% entropy reduction threshold on non-structured data.

\textbf{Algorithm 2 byte-match profiles.} The following table shows
Algorithm 2's raw byte-match frequency output on each dataset,
demonstrating that it independently discovers the correct record widths:

\begin{longtable}[]{@{}
  >{\raggedright\arraybackslash}p{(\columnwidth - 10\tabcolsep) * \real{0.0874}}
  >{\raggedright\arraybackslash}p{(\columnwidth - 10\tabcolsep) * \real{0.2524}}
  >{\raggedright\arraybackslash}p{(\columnwidth - 10\tabcolsep) * \real{0.1553}}
  >{\raggedright\arraybackslash}p{(\columnwidth - 10\tabcolsep) * \real{0.2039}}
  >{\raggedright\arraybackslash}p{(\columnwidth - 10\tabcolsep) * \real{0.1456}}
  >{\raggedright\arraybackslash}p{(\columnwidth - 10\tabcolsep) * \real{0.1553}}@{}}
\toprule\noalign{}
\begin{minipage}[b]{\linewidth}\raggedright
Dataset
\end{minipage} & \begin{minipage}[b]{\linewidth}\raggedright
Top byte-match lag (freq)
\end{minipage} & \begin{minipage}[b]{\linewidth}\raggedright
2nd lag (freq)
\end{minipage} & \begin{minipage}[b]{\linewidth}\raggedright
Candidates validated
\end{minipage} & \begin{minipage}[b]{\linewidth}\raggedright
Detected width
\end{minipage} & \begin{minipage}[b]{\linewidth}\raggedright
Matches Alg 1?
\end{minipage} \\
\midrule\noalign{}
\endhead
\bottomrule\noalign{}
\endlastfoot
sensor\_log & 3 (0.062) & 11 (0.062) & 3 & 8 & Yes \\
intel\_lab & 24 (0.365) & 48 (0.314) & 5 & 24 & Yes \\
mr & 2 (0.407) & 4 (0.384) & 5--6 & 2* & Yes \\
sao & 56 (0.076) & 28 (0.071) & 5 & 28 & Yes \\
NAS\-DAQ ITCH & 36 (0.843) & 8 (0.338) & 4 & 36† & No (Alg 1: 12) \\
\end{longtable}

*mr: detected on \textasciitilde30\% of blocks; remaining blocks return
\(w^* = 0\). †NAS\-DAQ ITCH: Algorithm 2 detects w=36 (true record width),
but the per-block compression comparison rejects it
(\(G = 1.00\times\)). Algorithm 1 detects w=12 (a divisor of 36, present
in the candidate set), also rejected by the compression comparison. The
strong byte-match peak at lag=36 (0.843) confirms that Algorithm 2 can
identify record widths outside Algorithm 1's candidate set.

For sensor\_log, the byte-match phase produces weak peaks at lags 3, 11,
19, 27 (period-8 offsets where specific byte positions happen to share
values); the correct width=8 is recovered via the fallback set and
entropy validation, not the byte-match mechanism.

\hypertarget{universality-gap-measurements}{%
\subsubsection{7.2 Universality Gap
Measurements}\label{universality-gap-measurements}}

\begin{longtable}[]{@{}
  >{\raggedright\arraybackslash}p{(\columnwidth - 14\tabcolsep) * \real{0.0763}}
  >{\raggedright\arraybackslash}p{(\columnwidth - 14\tabcolsep) * \real{0.1102}}
  >{\raggedright\arraybackslash}p{(\columnwidth - 14\tabcolsep) * \real{0.1610}}
  >{\raggedright\arraybackslash}p{(\columnwidth - 14\tabcolsep) * \real{0.2712}}
  >{\raggedright\arraybackslash}p{(\columnwidth - 14\tabcolsep) * \real{0.0763}}
  >{\raggedright\arraybackslash}p{(\columnwidth - 14\tabcolsep) * \real{0.1102}}
  >{\raggedright\arraybackslash}p{(\columnwidth - 14\tabcolsep) * \real{0.1102}}
  >{\raggedright\arraybackslash}p{(\columnwidth - 14\tabcolsep) * \real{0.0847}}@{}}
\toprule\noalign{}
\begin{minipage}[b]{\linewidth}\raggedright
Dataset
\end{minipage} & \begin{minipage}[b]{\linewidth}\raggedright
Block Entropy
\end{minipage} & \begin{minipage}[b]{\linewidth}\raggedright
\(R_{\text{ZSTD}}\)
\end{minipage} & \begin{minipage}[b]{\linewidth}\raggedright
\(R_{\text{shuffle+delta+ZSTD}}\)
\end{minipage} & \begin{minipage}[b]{\linewidth}\raggedright
Gap \(G\)
\end{minipage} & \begin{minipage}[b]{\linewidth}\raggedright
MSB Entropy
\end{minipage} & \begin{minipage}[b]{\linewidth}\raggedright
LSB Entropy
\end{minipage} & \begin{minipage}[b]{\linewidth}\raggedright
Gradient
\end{minipage} \\
\midrule\noalign{}
\endhead
\bottomrule\noalign{}
\endlastfoot
sensor\_log (synth.) & 5.5 bits & 4.25:1 & 173.5:1 & 40.8× & 0.00 & 5.00
& 5.00 \\
intel\_lab (real) & 6.1 bits & 3.58:1 & 4.26:1 & 1.19× & 0.17 & 6.07 &
5.90 \\
mr (real) & 2.7 bits & 2.72:1 & 2.76:1 & 1.02× & 0.47 & 5.05 & 4.58 \\
sao (real) & 7.4 bits & 1.25:1 & 1.26:1 & 1.01× & 3.19 & 6.08 & 2.89 \\
NAS\-DAQ ITCH (real) & 5.3 bits & 2.80:1 & 2.80:1 & 1.00× &
\textasciitilde6.0 & \textasciitilde7.5 & \textasciitilde1.5 \\
\end{longtable}

The gap factor \(G\) correlates with the MSB-to-LSB entropy gradient
when controlling for absolute MSB entropy (Section 4.4). The
relationship is not monotonic in the gradient alone: Intel Lab sensors
(gradient 5.90, \(G = 1.19\)) have a larger gradient than synthetic
sensor data (gradient 5.00, \(G = 40.8\)), but a smaller gap --- because
their MSB entropy (0.17 bits) is nonzero. The 40.8× gap is measured on
synthetic telemetry with near-ideal structure; the real-world results
(1.00×--1.19×) show the range.

\hypertarget{per-plane-entropy-profiles}{%
\subsubsection{7.3 Per-Plane Entropy
Profiles}\label{per-plane-entropy-profiles}}

\textbf{Sensor telemetry, synthetic (w=8):} 4 constant planes (entropy
\textless{} 0.01 after delta), 4 sparse planes (0.01--0.45, including
the carry-propagation plane at 0.45). Total correlation:
\textasciitilde17 bits per 64 KB block.

\textbf{Intel Lab sensors, real (w=24):} The three-class partition
appears but with higher entropies than the synthetic case. MSB planes of
float32 fields (positions 7, 11 within each 4-byte field) average
0.17--0.54 bits after delta, with LSB planes (positions 0, 1) average
2.3--6.1 bits. Mid-byte planes (positions 2, 6, 10) average 1.5--3.0
bits after delta. The partition is noisier than the synthetic case ---
real sensor data has irregular timestamps, missing readings, and 54
distinct sensor IDs that introduce more byte-level variation.

\textbf{MRI (w=2):} 1 constant-class MSB plane (entropy 0.15 after
delta, capturable on 31\% of blocks), 1 medium LSB plane (entropy 4.22
after delta). Width=2 is auto-detected on \textasciitilde30\% of blocks
(uniform-tissue regions); on the remaining blocks, the 20\% entropy
reduction threshold is not met and the algorithm correctly falls back to
plain ZSTD.

\textbf{SAO (w=28):} 7 low-entropy MSB planes (position 3 mod 4, entropy
\textasciitilde3.2), 7 medium planes (position 2 mod 4, entropy
\textasciitilde5.3), 14 high-entropy planes (positions 0,1 mod 4,
entropy \textasciitilde6.1). No planes are capturable by simple
generators --- the per-field entropy is too high.

\hypertarget{impact-on-compression-benchmarks}{%
\subsubsection{7.4 Impact on Compression
Benchmarks}\label{impact-on-compression-benchmarks}}

On the Silesia Corpus (12 files, 202 MB total), the BPD implementation
achieves 2.90:1 aggregate ratio versus ZSTD per-block at 2.90:1 --- no
regression on non-structured data. Per-file results for files where the
width detection algorithm activates:

\begin{longtable}[]{@{}
  >{\raggedright\arraybackslash}p{(\columnwidth - 10\tabcolsep) * \real{0.0750}}
  >{\raggedright\arraybackslash}p{(\columnwidth - 10\tabcolsep) * \real{0.0750}}
  >{\raggedright\arraybackslash}p{(\columnwidth - 10\tabcolsep) * \real{0.2625}}
  >{\raggedright\arraybackslash}p{(\columnwidth - 10\tabcolsep) * \real{0.3000}}
  >{\raggedright\arraybackslash}p{(\columnwidth - 10\tabcolsep) * \real{0.1000}}
  >{\raggedright\arraybackslash}p{(\columnwidth - 10\tabcolsep) * \real{0.1875}}@{}}
\toprule\noalign{}
\begin{minipage}[b]{\linewidth}\raggedright
File
\end{minipage} & \begin{minipage}[b]{\linewidth}\raggedright
Type
\end{minipage} & \begin{minipage}[b]{\linewidth}\raggedright
\(R_{\text{BPD}}\)
\end{minipage} & \begin{minipage}[b]{\linewidth}\raggedright
\(R_{\text{ZSTD-block}}\)
\end{minipage} & \begin{minipage}[b]{\linewidth}\raggedright
Change
\end{minipage} & \begin{minipage}[b]{\linewidth}\raggedright
Detected Width
\end{minipage} \\
\midrule\noalign{}
\endhead
\bottomrule\noalign{}
\endlastfoot
mr & MRI & 2.76:1 & 2.72:1 & \textbf{+1.5\%} & 2 \\
sao & Star catalog & 1.25:1 & 1.25:1 & 0\%* & 28 \\
\end{longtable}

*SAO's aggregate ratio ties ZSTD-block because per-block strategy
selection chooses plain ZSTD on blocks where shuffle doesn't help. The
per-block shuffle improvement averages \textasciitilde2.4\% on blocks
where it activates, but not all 111 blocks benefit, so the aggregate is
a tie.

On text files (dickens 2.42:1, webster 3.02:1, reymont 3.07:1), the BPD
implementation ties ZSTD within 0.5\% --- width detection correctly
returns \(w^* = 0\) and the data passes through to plain ZSTD. The
aggregate 2.90:1 includes these files where width detection does
nothing, demonstrating that the algorithm adds no regression on
non-structured data.

On the Intel Berkeley Lab real-world sensor dataset {[}16{]}, the BPD
implementation achieves 4.26:1 versus ZSTD per-block at 3.58:1
(\textbf{+19.0\%} improvement, universality gap \(G = 1.19\)). The
automatically detected width of 24 matches the 24-byte record structure.

\hypertarget{computational-overhead}{%
\subsubsection{7.5 Computational
Overhead}\label{computational-overhead}}

Width detection adds computational cost to the compression path only;
decompression is unaffected since the detected width is stored in the
block header.

\begin{longtable}[]{@{}
  >{\raggedright\arraybackslash}p{(\columnwidth - 6\tabcolsep) * \real{0.2121}}
  >{\raggedright\arraybackslash}p{(\columnwidth - 6\tabcolsep) * \real{0.2576}}
  >{\raggedright\arraybackslash}p{(\columnwidth - 6\tabcolsep) * \real{0.2424}}
  >{\raggedright\arraybackslash}p{(\columnwidth - 6\tabcolsep) * \real{0.2879}}@{}}
\toprule\noalign{}
\begin{minipage}[b]{\linewidth}\raggedright
Configuration
\end{minipage} & \begin{minipage}[b]{\linewidth}\raggedright
Candidate Widths
\end{minipage} & \begin{minipage}[b]{\linewidth}\raggedright
Shuffles/Block
\end{minipage} & \begin{minipage}[b]{\linewidth}\raggedright
Detection Overhead
\end{minipage} \\
\midrule\noalign{}
\endhead
\bottomrule\noalign{}
\endlastfoot
Original (v0.4) & \{2, 4, 8\} & 3 & Baseline \\
Algorithm 1 (v0.5) & 15 fixed widths & 15 & +12\% \\
Algorithm 2 (v0.6) & \textasciitilde5--8 adaptive & 5--8 & +4--7\% \\
\end{longtable}

Each candidate width requires one shuffle operation (\(O(n)\)) and one
entropy computation (\(O(n)\)), for a total cost of \(O(|W| \cdot n)\)
per block. On 64 KB blocks, this amounts to \textasciitilde100--200
\(\mu\)s per candidate width. Algorithm 2 adds a sampled byte-match
phase (\(O(L \cdot 4096) \approx 260K\) comparisons, negligible versus
shuffle cost) but reduces the number of validated candidates from 15 to
typically 5--8, yielding a net reduction in overhead.

Algorithm 2 removes the fixed candidate set limitation: it detects
arbitrary widths up to \(L = 64\) without a predefined set, which means
widths like 28 (SAO catalog) and 24 (Intel Lab) are discovered
automatically rather than requiring manual inclusion. Both algorithms
produce identical width detection results on four of five tested
datasets; on NAS\-DAQ ITCH, Algorithm 2 detects the true width=36 while
Algorithm 1 detects w=12, with both rejected by the per-block
compression comparison (Section 7.1).

\begin{center}\rule{0.5\linewidth}{0.5pt}\end{center}

\hypertarget{block-size-sensitivity}{%
\subsubsection{7.6 Block Size
Sensitivity}\label{block-size-sensitivity}}

We measure the universality gap at block sizes from 32 KB to 256 KB:

\begin{longtable}[]{@{}lllll@{}}
\toprule\noalign{}
Dataset & 32 KB & 64 KB & 128 KB & 256 KB \\
\midrule\noalign{}
\endhead
\bottomrule\noalign{}
\endlastfoot
sensor\_log & 32.6× & 40.8× & 46.0× & 54.6× \\
intel\_lab & 1.18× & 1.19× & 1.19× & 1.13× \\
\end{longtable}

For highly regular data (sensor\_log), the gap increases with block size
because larger planes provide more samples for entropy estimation and
better ZSTD dictionary utilization. For real-world sensor data
(intel\_lab), the gap is stable across all tested block sizes
(1.13×--1.19×), confirming that the results are not an artifact of the
64 KB block size choice.

\begin{center}\rule{0.5\linewidth}{0.5pt}\end{center}

\hypertarget{discussion}{%
\subsection{8. Discussion}\label{discussion}}

\hypertarget{limitations}{%
\subsubsection{8.1 Limitations}\label{limitations}}

\textbf{Width detection has a bounded search range.} Algorithm 2
supports arbitrary widths up to the maximum lag \(L = 64\). Records
wider than 64 bytes would require increasing \(L\), which linearly
increases the byte-match phase cost. The \(L = 64\) bound does not cover
the 90-byte enterprise packet captures {[}19{]} cited as future targets;
increasing \(L\) to 128 would address this with proportional cost
increase in the byte-match phase. Note that reliable detection at lag=90
on 64 KB blocks yields only \textasciitilde4090 samples (barely
sufficient), so larger record widths may also require larger blocks. In
practice, the vast majority of binary record formats use widths under 64
bytes.

\textbf{The universality gap is data-dependent.} Our five-point
characterization spans from no improvement (NAS\-DAQ ITCH, \(G = 1.00\))
to 40.8× (synthetic sensor telemetry). Promising future targets for the
2×--10× range include USGS LiDAR LAS point clouds (30-byte records)
{[}18{]} and enterprise packet captures (90-byte records) {[}19{]}.
Standard compression benchmarks (Canterbury, LTCB, lzbench) lack
structured binary data entirely --- developing a standard binary
benchmark corpus is a community need.

\textbf{Entropy reduction does not guarantee compression improvement.}
NAS\-DAQ ITCH demonstrates this clearly: Algorithm 2 detects the true
width=36 with 48\% entropy reduction, yet the universality gap is 1.00×
because ZSTD's LZ matching already captures the interleaved patterns.
Order-0 Shannon entropy (which the detection threshold uses) does not
account for the higher-order redundancy that ZSTD exploits. The
two-stage design (detection + per-block compression comparison) is
essential to prevent regressions in such cases.

\textbf{The threshold \(\tau = 0.80\) is a sharp optimum.} Section 3.5
shows that \(\tau = 0.80\) is the most conservative threshold achieving
perfect detection on all structured datasets, with more margin than the
only alternative (\(\tau = 0.90\)) that also achieves 5/5. The two-stage
architecture (detection + per-block comparison) mitigates this
fragility: false positives from a relaxed threshold are caught by the
compression comparison, which only uses shuffle when it produces shorter
output. This defense-in-depth means the system never regresses even with
suboptimal \(\tau\).

\textbf{Per-plane generative capture has limited benefit.} While the
three-class partition identifies constant planes that could be stored as
generator formulas, ZSTD already compresses these planes extremely
efficiently within the joint compression frame. The overhead of
per-plane metadata (2--10 bytes per plane) often exceeds the marginal
savings.

\hypertarget{relation-to-ica-and-mdl}{%
\subsubsection{8.2 Relation to ICA and
MDL}\label{relation-to-ica-and-mdl}}

Byte-plane decomposition can be interpreted as Independent Component
Analysis (ICA) over finite alphabets {[}11{]}, where the ``mixing
matrix'' is the known byte-interleaving pattern of struct serialization
and the ``demixing'' is the fixed shuffle permutation. Unlike general
ICA, no estimation is needed --- only the period (width) must be
determined.

The width selection and per-block strategy selection follow the Minimum
Description Length (MDL) principle {[}12{]}: each candidate strategy
defines a model class, and the strategy with the shortest two-part code
(model description + data given model) is chosen. This is a discrete
approximation of the Kolmogorov structure function {[}13{]}.

\hypertarget{the-carry-propagation-phenomenon}{%
\subsubsection{8.3 The Carry-Propagation
Phenomenon}\label{the-carry-propagation-phenomenon}}

An interesting inter-plane dependency emerges in the sensor telemetry
data: plane 1 (byte position 1 of the 4-byte timestamp) exhibits a
bimodal distribution determined by the overflow frequency of plane 0.
The carry frequency is exactly
\((\delta \bmod 256) / 256 = 232/256 \approx 0.906\), matching the
observed 90.6\% distribution. This cross-plane dependency is predictable
from the field-level model (a 4-byte integer with constant increment
1000) but cannot be captured by any single-plane pattern detector. The
optimal generative model for such data operates at the \textbf{field
level}, not the byte-plane level --- a finding that suggests field-level
generators as a direction for future work.

\begin{center}\rule{0.5\linewidth}{0.5pt}\end{center}

\hypertarget{related-work}{%
\subsection{9. Related Work}\label{related-work}}

\textbf{Bitshuffle} {[}4{]} (Masui et al., 2015) provides an intuitive
argument for byte-plane decomposition's effectiveness on numerical data
but offers no formal analysis or auto-detection algorithm.

\textbf{TDT} {[}7{]} (arXiv:2506.18062, 2025) goes beyond fixed-width
shuffling to cluster byte positions by entropy features. Their key
finding --- that the benefit operates at the level of higher-order
entropy --- is consistent with our observations. However, TDT does not
auto-detect the record width and requires more complex clustering
machinery.

\textbf{ZipNN} {[}8{]} (Hershcovitch et al., 2024) demonstrates
byte-plane decomposition for AI model compression, achieving significant
savings by separating exponent bytes (highly skewed) from mantissa bytes
(near-uniform). This is a domain-specific application of the same
principle we formalize here.

\textbf{Blosc bytedelta} {[}5{]} (Alted, 2023) chains shuffle and delta
encoding, reporting \textasciitilde15\% improvement over shuffle alone
on climate data. Our delta encoding follows the same approach.

\textbf{Column-oriented databases} {[}14, 15{]} achieve 5--10× better
compression than row-oriented storage by type-homogeneous per-column
encoding. Byte-plane decomposition achieves a similar effect at the byte
level without schema knowledge.

\begin{center}\rule{0.5\linewidth}{0.5pt}\end{center}

\hypertarget{conclusion}{%
\subsection{10. Conclusion}\label{conclusion}}

We have presented two algorithms for automatically detecting the optimal
byte-plane decomposition width via entropy minimization, and the first
systematic characterization of the universality gap between
structure-aware and structure-blind lossless compression. Algorithm 1
evaluates a fixed candidate set of 15 widths; Algorithm 2 uses sampled
byte-match frequency to detect periodic byte repetitions at arbitrary
lags up to 64, reducing the number of validated candidates to 5--8 while
supporting any record width without a predefined set. Algorithm 2
successfully detected the 36-byte record width in NAS\-DAQ ITCH data
(outside Algorithm 1's candidate set), confirming its ability to
discover arbitrary widths. Both algorithms require a two-stage design:
entropy-based detection followed by per-block compression comparison,
since entropy reduction alone does not guarantee compression improvement
(NAS\-DAQ: 48\% entropy reduction, \(G = 1.00\times\)).

Our five-point empirical curve --- with gaps ranging from 1.00× (NAS\-DAQ
financial data) to 40.8× (synthetic sensor telemetry) --- demonstrates
that the gap is predicted by the MSB-to-LSB entropy gradient and the
absolute MSB entropy within multi-byte fields, providing a practical
criterion for when byte-plane decomposition will be beneficial. The
three-class plane partition (constant/sparse/medium) that emerges after
decomposition follows from elementary properties of bounded integer
representations and provides a framework for understanding which planes
benefit from decomposition.

These results suggest that the ``price of universality'' in lossless
compression is not merely asymptotic overhead but can represent
orders-of-magnitude performance gaps on highly regular structured binary
data --- gaps that simple, schema-free preprocessing can largely close.
The schema-free nature of the width detection algorithm is its key
practical advantage: unlike HDF5, Blosc, and Parquet, it requires no
user configuration or data format knowledge, making it applicable as a
transparent preprocessing layer in any compression pipeline.

Future work includes validation on binary datasets with estimated gaps
in the 2×--10× range (USGS LiDAR point clouds {[}18{]}, enterprise
packet captures {[}19{]}) and formal bounds on the universality gap
using LZ optimality theory {[}9{]}.

\begin{center}\rule{0.5\linewidth}{0.5pt}\end{center}

\hypertarget{acknowledgments}{%
\subsection{Acknowledgments}\label{acknowledgments}}

This work was developed with substantial assistance from AI tools. The
author conceived the byte-plane decomposition approach, the
auto-detection algorithm, and the universality gap framework. Claude
(Anthropic) was used extensively for literature research, theoretical
analysis, experimental design feedback, and iterative manuscript review.
All experimental results were produced by the author's implementation
and independently verified via automated reproducibility scripts.

\begin{center}\rule{0.5\linewidth}{0.5pt}\end{center}

\hypertarget{references}{%
\subsection{References}\label{references}}

{[}1{]} Y. Collet, ``Zstandard compression,'' RFC 8478, 2018.

{[}2{]} J. Ziv and A. Lempel, ``A universal algorithm for sequential
data compression,'' IEEE Trans. Information Theory, vol.~23, no. 3,
pp.~337--343, 1977.

{[}3{]} The HDF Group, ``HDF5 shuffle filter,''
\url{https://docs.hdfgroup.org/hdf5/}.

{[}4{]} K. Masui et al., ``Efficient compression of radio
interferometric data,'' Astronomy and Computing, vol.~12, pp.~181--190,
2015.

{[}5{]} F. Alted, ``Bytedelta: enhance your compression toolset,'' Blosc
Blog, 2023,
\url{https://blosc.org/posts/bytedelta-enhance-compression-toolset/}.

{[}6{]} Apache Software Foundation, ``Apache Parquet format
specification: BYTE\_STREAM\_SPLIT encoding,''
\url{https://parquet.apache.org/}.

{[}7{]} TDT: Typed Data Transformation, arXiv:2506.18062, 2025.

{[}8{]} M. Hershcovitch et al., ``ZipNN: lossless compression for AI
models,'' arXiv:2411.05239, 2024.

{[}9{]} P. Ferragina, I. Nitto, and R. Venturini, ``On the
bit-complexity of Lempel-Ziv compression,'' SIAM J. Computing, vol.~42,
no. 4, pp.~1521--1541, 2013.

{[}10{]} Silesia Compression Corpus,
\url{https://sun.aei.polsl.pl/~sdeor/index.php?page=silesia}.

{[}11{]} A. Painsky, S. Rosset, and M. Feder, ``Generalized independent
component analysis over finite alphabets,'' IEEE Trans. Information
Theory, vol.~62, no. 2, pp.~1038--1053, 2016.

{[}12{]} P. Grünwald, The Minimum Description Length Principle, MIT
Press, 2007.

{[}13{]} N. Vereshchagin and P. Vitányi, ``Kolmogorov's structure
functions and model selection,'' IEEE Trans. Information Theory,
vol.~50, no. 12, pp.~3265--3290, 2004.

{[}14{]} D. Abadi, S. Madden, and M. Ferreira, ``Integrating compression
and execution in column-oriented database systems,'' Proc. SIGMOD,
pp.~671--682, 2006.

{[}15{]} D. Abadi, P. Boncz, and S. Harizopoulos, ``Column-oriented
database systems,'' Proc. VLDB Endowment, vol.~2, no. 2, pp.~1664--1665,
2009.

{[}16{]} P. Bodik et al., ``Intel Berkeley Research Lab sensor data,''
2004, \url{http://db.csail.mit.edu/labdata/labdata.html}.

{[}17{]} NAS\-DAQ, ``ITCH 5.0 total view data,''
\url{https://emi.nasdaq.com/ITCH/Nasdaq\%20ITCH/}.

{[}18{]} USGS, ``3D Elevation Program (3DEP) LiDAR data,''
\url{https://apps.nationalmap.gov/downloader/}.

{[}19{]} R. Pang et al., ``The devil and packet trace anonymization,''
ACM Computer Communication Review, vol.~36, no. 1, pp.~29--38, 2006.

\end{document}
